\documentclass[11pt]{amsart}
\usepackage{amssymb,mathtools,hyperref,booktabs}
\usepackage[margin=1.1in]{geometry}
% compile with: latexmk -pdf note.tex

\newtheorem{theorem}{Theorem}
\newtheorem{lemma}[theorem]{Lemma}
\newtheorem{proposition}[theorem]{Proposition}
\newtheorem{conjecture}[theorem]{Conjecture}
\theoremstyle{remark}
\newtheorem{remark}[theorem]{Remark}
\DeclareMathOperator{\adj}{adj}
\DeclareMathOperator{\per}{per}
\DeclareMathOperator{\codim}{codim}
\DeclareMathOperator{\len}{length}
\DeclareMathOperator{\Ann}{Ann}
\newcommand{\bdc}{\overline{\mathrm{dc}}}

\title{A remark on OpenAI family 108: the cubic constant for $\bdc(\per_m)$}
\author{[AUTHOR NAME]}
\date{October 9, 2026}

\begin{document}
\begin{abstract}
OpenAI's manuscript \emph{A cubic lower bound for border determinantal
complexity of the permanent} (family~108 of \texttt{github.com/openai/math},
commit \texttt{adc7f124}) proves $\bdc(\per_m)\ge m^3/(5\,529\,600\,e)$ for
$m\ge1408$. We record three facts.
(1) A short lemma (``box lemma'') improves the manuscript's length bound
$\len Z\ge\binom{q+2}2$ to $\len Z\ge(q+1)^2$ and removes the hypothesis
$q\le r-1$. Together with re-chosen parameters this gives
$\bdc(\per_m)\ge m^3/(125\,500\,e)$ for all $m\ge1408$, and
$m^3/((118\,800+o(1))e)$ asymptotically: $44\times$ resp.\ $46.5\times$
better. Parameter choice alone accounts for about $21\times$/$23\times$;
the box lemma for the remaining factor $2.1\times$/$2.0\times$.
(2) The gradient locus contains explicit block-triangular families, valid in
every characteristic, whose codimension is asymptotically
$(\tfrac12-\beta^2)k^2$ with $\beta=r/k$. This is smaller than
$(2q+2)^2$, the codimension of the rank family. So the latter is not the true
codimension, and two natural strengthening lemmas are false.
(3) Consequently no proof inside the manuscript's pipeline (selector, smooth
linear section, smooth-hypersurface theorem) can give a uniform constant
better than $1/((33\,065+o(1))e)$, i.e.\ more than $167\times$ the stated one.
The construction and all main arguments are OpenAI's.
\end{abstract}
\maketitle

\section{Setting}
We use the notation of the manuscript. $K$ is algebraically closed,
$\theta_1,\dots,\theta_k\in K$ are distinct, $D=\mathrm{diag}(\theta_i)$,
$0\le q\le (k-2)/2$, $r=k-2q\ge2$, $c=r-1$, and
\[
G_D(X)=[s^qt^q]\det(X+sI+tD),\qquad
\Gamma_D=\{X:\nabla G_D(X)=0\}\subset\mathrm{Mat}_k(K).
\]
For fixed $X$ let $M=sI+tD+uX$, $F_X=\det M$, $C_X=V(F_X)\subset\mathbb P^2$,
and let $Z=Z_X\subset C_X$ be the finite scheme cut out by the entries of
$\adj M$ (\texttt{adjugate.tex} L123--137). Let $I_Z\subset S=K[s,t,u]$ be its
saturated homogeneous ideal.

The manuscript's chain is: $\nabla G_D(X)=0\iff[s^qt^qu^c]\adj M=0$
(\texttt{adjugate.tex} L261--270, no hypothesis on $q$);
$\len Z\ge T:=\binom{q+2}2$ (Prop.~\texttt{adjugate-length}, L250--309, which
uses $q\le c$ at L291); $\len Z\le\delta(C_X)$
(Prop.~\texttt{conductor-length}); hence
$\dim\Gamma_D\le k^2+k-T$ (\texttt{dimension.tex} L10--27, L353--366).
With $\Delta=T-k$, characteristic transfer, a smooth linear section of
dimension $d$ with $2\le d<k^2$, $d-1<\Delta$, and the selector
$m\ge11k-2q$ (\texttt{transfer.tex} L218--309, \texttt{selector.tex} L55--58),
Theorem~\texttt{smooth-border} gives
\begin{equation}\label{eq:chain}
\bdc(\per_m)\ \ge\ \frac{(r-1)(d-1)}{4e}.
\end{equation}
The manuscript takes $q=\lfloor(k-1)/3\rfloor$, $d=\lfloor k^2/100\rfloor$,
$k=\lfloor m/11\rfloor$ (\texttt{consequences.tex} L8--38).

\sloppy
\paragraph{Related work.} Two follow-ups bound other quantities and leave the
cubic constant $5\,529\,600$ unchanged. Guo~\cite{Guo} (v2, October 8, 2026)
moves the smooth slice of family~108 to circuits with division and obtains
$L_{\mathrm{div}}(\per_n)\ge(n^2/86400)\log_2(n/48)$ ($n\ge1408$), assuming the
manuscript's Prop.~7.3; from the Lean declarations alone he obtains
$(n^2/119790)\log_2(n/44)$ ($n\ge1936$). Xie~\cite{Xie} (October 2, 2026)
proves $\bdc(\per_n)=\Omega(n^2\log n)$, explicitly
$\ge n^2(\log_2n-5)/(64e)$ for $n\ge64$.

\paragraph{Paper versus Lean.} Guo also points out that the Lean slice size
$D=\lfloor\binom{q+2}2/18\rfloor$ differs from Prop.~7.3. We confirm this. The
Lean development (\texttt{Main.lean}, \texttt{critical\_dimension\_effective})
proves only $\dim\Gamma_D\le k^2-\lfloor T/18\rfloor$, where the paper's
Prop.~6.1 has $k^2+k-T$. It reaches the same constant $5\,529\,600$ through
tighter auxiliary estimates ($r-1\ge0.33k$, $D-1\ge k^2/330$, $m\le11.1k$).
Our baseline and all our constants follow the paper: Prop.~6.1, $d\le\Delta$,
and $5\,529\,600=12^3\cdot3200$ from \texttt{consequences.tex}. The number is
the same in both sources, so the ratios reported below are unaffected. However,
our constants rest on the paper's Prop.~6.1 and are not Lean-certified. If one
uses only the Lean dimension bound, re-chosen parameters give
$c^{-1}\approx4.2\cdot10^6e$ ($1.3\times$).

\section{The box lemma}
\begin{lemma}[Box lemma]\label{lem:box}
Let $0\le q\le (k-1)/2$ and $X\in\Gamma_D$. Then $\len Z_X\ge(q+1)^2$.
\end{lemma}

\begin{proof}
Put $\Lambda=[s^qt^qu^c]$, a linear functional on $S_{k-1}$.

\emph{Step 1.} By Lemma~\texttt{adjugate-generation} (\texttt{adjugate.tex}
L139--248, no hypothesis on $q$), every form of degree $k-1$ in $I_Z$ is a
constant linear combination of entries of $\adj M$. These entries are killed
by $\Lambda$ because $X\in\Gamma_D$. Let $P\in(I_Z)_d$ with $d\le k-1$. For every
monomial $Q$ of degree $k-1-d$ we have $PQ\in(I_Z)_{k-1}$, so $\Lambda(PQ)=0$.
In other words, no monomial dividing $s^qt^qu^c$ occurs in $P$. Thus
\[
(I_Z)_d\subset(s^{q+1},t^{q+1},u^{c+1})_d\qquad(d\le k-1).
\]
This is coefficient extraction, not differentiation, so it holds in every
characteristic, including~$2$. It is the computation of \texttt{adjugate.tex}
L283--297, applied to all degrees and not only to degree~$q$.

\emph{Step 2.} Let $B=(s^{q+1},t^{q+1},u)$. For $d\le 2q\le k-1$, Step~1 gives
$(I_Z+(u))_d\subset B_d$. For $d\ge2q+1$, every monomial of degree $d$ in $s,t$
is divisible by $s^{q+1}$ or $t^{q+1}$, so $B_d=S_d$. Hence
$I_Z+(u)\subset B$ in every degree.

\emph{Step 3.} $Z$ does not meet the line $u=0$. The points of $C_X$ on it are
$P_i=[-\theta_i:1:0]$, and there the diagonal adjugate entry
$\prod_{j\ne i}(\theta_j-\theta_i)$ is nonzero (\texttt{adjugate.tex} L81--87).
So $u$ lies in no associated prime of $S/I_Z$ and is a nonzerodivisor. Then
$H_{S/(I_Z+u)}(d)=H_{S/I_Z}(d)-H_{S/I_Z}(d-1)$, and summing gives
\[
\len Z=\dim_K S/(I_Z+(u))\ \ge\ \dim_K S/B=(q+1)^2. \qedhere
\]
\end{proof}

\begin{remark}
The manuscript's argument extracts only the injectivity of
$S_q\to H^0(\mathcal O_Z(q))$, which gives $\dim S_q=\binom{q+2}2$. Step~3
turns the whole Hilbert function of the Artinian reduction into length. The
lemma is short; we regard it as a small improvement, not as new technique.
\end{remark}

Replacing $T$ by $(q+1)^2$ in \texttt{dimension.tex} gives
$\dim\Gamma_D\le k^2+k-(q+1)^2$ for every $0\le q\le (k-2)/2$. The Severi bound,
the conductor bound, and the fibre bound do not involve $q$.

\section{Constants}
With $\Delta=(q+1)^2-k$ we take $d=\Delta$ (so $d-1<\Delta$), which needs
$2\le\Delta<k^2$, $r\ge2$, and $11k-2q\le m$. We maximise $(r-1)(d-1)$ in
\eqref{eq:chain}. Write $c^{-1}=4m^3/((r-1)(d-1))$ in units of $e$.

\begin{theorem}\label{thm:new}
For every integer $m\ge1408$, $\bdc(\per_m)\ge \dfrac{m^3}{125\,500\,e}$.
Moreover $\bdc(\per_m)\ge \dfrac{m^3}{(118\,800+o(1))\,e}$ as $m\to\infty$.
The same bounds hold for $\mathrm{dc}(\per_m)$.
\end{theorem}

\begin{proof}[Proof sketch]
\emph{Asymptotics.} Put $q=\gamma k$. Then $m\approx(11-2\gamma)k$ and
$(r-1)(d-1)\approx(1-2\gamma)\gamma^2k^3$. The function
$(1-2\gamma)\gamma^2/(11-2\gamma)^3$ is maximised at $\gamma=11/32$, with value
$1/29\,700$. So $c^{-1}\to4\cdot29\,700=118\,800$.

\emph{Tail.} Take $k=\lfloor 32(m-2)/330\rfloor$ and $q=\lfloor 11k/32\rfloor$.
Then $11k-2q\le m$, $r-1\ge 10k/32-1$, and $d-1\ge(11k/32)^2-k-1$. Substituting
$k\ge 32(m-2)/330-1$ gives an explicit rational function $g(m)\ge c^{-1}(m)$.
The numerator of $g'$ has no real root above $65$ (exact computation), so
$g$ decreases. Also $g(200\,000)=118\,893$.

\emph{Range $1408\le m\le200\,000$.} For each $m$ we search $(k,q)$ near the
optimum. The largest value of $c^{-1}(m)$ found is $125\,422.2$, at $m=1419$.
See \texttt{scripts/independent/uniform\_indep.py}.
\end{proof}

\begin{table}[h]
\begin{tabular}{lrrr}
\toprule
 & uniform, $m\ge1408$ & $m=5000$ & $m\to\infty$\\
\midrule
stated (\texttt{introduction.tex} L48) & $5\,529\,600$ & $5\,529\,600$ & $5\,529\,600$\\
manuscript's lemma, re-chosen parameters & $267\,955$ & $245\,793$ & $238\,328$\\
box lemma, re-chosen parameters & $125\,422$ & $120\,334$ & $118\,800$\\
\midrule
gain from parameters only & $20.6\times$ & $22.5\times$ & $23.2\times$\\
gain from box lemma & $2.14\times$ & $2.04\times$ & $2.006\times$\\
total & $44.1\times$ & $46.0\times$ & $46.5\times$\\
\bottomrule
\end{tabular}
\caption{Values of $c^{-1}/e$. ``Re-chosen parameters'' means $q\le(k-1)/3$,
$d=\Delta$, and $k$ maximal subject to $11k-2q\le m$. Of the box-lemma factor
$2.006$, the change $\binom{q+2}2\to(q+1)^2$ contributes $2.000$, and dropping
$q\le c$ contributes $1.003$.}
\end{table}

Most of the loss in the stated constant comes from the manuscript's
deliberately crude estimates $d\approx k^2/100$ against $\Delta\approx k^2/18$,
$k\ge m/12$, and $r-1\ge k/4$ against $k/3$. These alone are worth
$\approx23\times$.

\section{The true codimension and the limit of the method}
\begin{proposition}[Block families]\label{prop:F}
The partials of $G_D$ are $K$-linear combinations of the
\emph{almost-principal} $(r-1)$-minors of $X$, that is, minors with row set
$U\cup\{a\}$ and column set $U\cup\{b\}$. Let $X$ be block upper triangular
for some ordered partition of the coordinates into $m'$ blocks, and let the
diagonal blocks have ranks $\rho_i$ with $\sum\rho_i+m'\le r-1$. Then every
such minor vanishes, so $X\in\Gamma_D$. This holds for every field, every
characteristic, and every $\theta$.
\end{proposition}

\begin{proof}
$\partial G_D/\partial X_{ij}=[s^qt^q]\adj(X+sI+tD)_{ji}$. Expanding along the
diagonal $(s+t\theta_l)$ entries gives a sum over $2q$-sets $S$ of a polynomial
in $\theta$ times the minor of $X$ with rows $[k]\setminus(S\cup\{i\})$ and
columns $[k]\setminus(S\cup\{j\})$. For the vanishing, apply the generalized
Laplace expansion along the row blocks:
$\det X[I,J]=\pm\sum\prod_p\det X[I_p,C_p]$, where $C_p$ consists of columns in
blocks $\ge p$. A term is nonzero only if $C_p$ contains at most $\rho_p$
columns of block $p$. Process the blocks in order. Since $|I_p|$ and $|J_p|$
differ in at most two blocks and by at most one, at most $m'-1$ columns are
assigned off the diagonal. So $\sum_p\rho_p\ge (r-1)-(m'-1)$, which
contradicts the hypothesis.
\end{proof}

Take one block of size $n_1\approx2\beta k$ and rank $r-1-m'$, and let the
remaining $m'-1\to\infty$ blocks have rank $0$ and equal sizes. The family has
codimension $(\tfrac12-\beta^2+o(1))k^2$, where $\beta=r/k$. The case $m'=1$
is the rank-$\le r-2$ family, of codimension $(2q+2)^2\approx(1-\beta)^2k^2$.
Since $(1-\beta)^2-(\tfrac12-\beta^2)=2(\beta-\tfrac12)^2$, the block families
are strictly larger for $\beta\ne\tfrac12$. Two explicit instances, each
verified by exact computation over $\mathbb Q$ and $\mathbb F_{2^{20}}$:
\begin{itemize}
\item $k=9$, $q=3$: $X=\begin{psmallmatrix}0&Y\\0&0\end{psmallmatrix}$ with
$4+5$ blocks has codimension $61<64=(2q+2)^2$.
\item $k=22$, $q=7$: blocks $(16,6)$, ranks $(5,0)$, codimension $253<256$.
\end{itemize}
Hence a lemma of the form $\codim\Gamma_D\ge(2q+2)^2-O(1)$ is false. The
natural intermediate lemma $\dim\{X:\len Z_X\ge\ell\}\le k^2+k-2\ell$ is also
false. For $k=16$ and $X=\begin{psmallmatrix}0&Y\\0&0\end{psmallmatrix}$ with
$8\times8$ blocks, every adjugate entry has order $\ge k-2$ at $[0:0:1]$.
So $\len Z_X\ge\binom{15}2=105$ for every $Y$, and the local length is exactly
$105$. This family has dimension $64>62$.

\begin{proposition}[Ceiling]\label{prop:ceiling}
Consider any proof that applies Theorem~\texttt{smooth-border} to a
$d$-dimensional linear section of $G_D$, for some diagonal $D$, and embeds it
in $\per_m$ by the selector ($m\ge11k-2q$). Then $d\le\codim\Gamma_D$, by the
projective dimension theorem and Euler's identity over $\mathbb C$.
Proposition~\ref{prop:F} therefore gives
\[
\frac{(r-1)(d-1)}{4e}\ \le\ \frac{(1+o(1))\,m^3}{C_{\rm ceil}\,e},\qquad
C_{\rm ceil}=\min_{0<\beta<1}\frac{4(10+\beta)^3}{\beta\min\{(1-\beta)^2,\tfrac12-\beta^2\}}
=33\,065.4\ldots,
\]
attained at $\beta=(\sqrt{601}-1)/60\approx0.392$. So no uniform constant from
this pipeline beats $1/((33\,065+o(1))e)$, a factor $167\times$ over the stated
constant. The box lemma is within $3.59\times$ of this ceiling.
\end{proposition}

At finite $m$ the per-$m$ ceiling is slightly weaker: $30\,872e$ at $m=1408$
($179\times$) and $31\,858e$ at $m=5000$ ($174\times$). The ceiling depends on
the premises stated. It does not cover initial forms at non-zero points,
other coefficients $[s^at^b]$, or other selectors.

The true codimension is therefore in
$\bigl[(q+1)^2-k,\ \min\{(2q+2)^2,\ \text{block-family codim}\}\bigr]$,
asymptotically $[\tfrac14(1-\beta)^2,\ \tfrac12-\beta^2]\,k^2$.

\begin{conjecture}[A$^*$, open]
For $\beta=r/k$ fixed in $(\tfrac14,\tfrac12)$,
$\codim\Gamma_D\ge(\tfrac12-\beta^2-o(1))k^2$ for generic $\theta$; that is,
the block families of Proposition~\ref{prop:F} are asymptotically extremal.
If true, the pipeline would reach $c^{-1}\to33\,065e$.
\end{conjecture}
We have not checked this conjecture beyond small cases, and we do not know
whether families with several ranked blocks do better.

\section{Computational certificates}
All scripts are in \texttt{scripts/}, with outputs in
\texttt{scripts/expected\_output/}. \texttt{box\_indep.sage} checks steps 1--3
of Lemma~\ref{lem:box} for $k\le9$ over $\mathbb Q$, $\mathbb F_{1000003}$ and
$\mathbb F_{2^{10}}$, including cases with $q>c$. \texttt{grad\_check.sage}
checks the counterexamples. It uses the characteristic-free identity
$\adj(sI-B)=\sum_j s^{k-1-j}\sum_{i\le j}c_iB^{j-i}$, which it cross-checks
against brute-force differentiation. \texttt{minor\_span.sage} checks the
almost-principal-minor span for $k\le7$. \texttt{lemmaB\_check.sage} computes
the local length $105$ exactly over $\mathbb Q$.

\begin{thebibliography}{9}
\bibitem{OAI} OpenAI, \emph{A cubic lower bound for border determinantal
complexity of the permanent}, September 24, 2026,
\texttt{github.com/openai/math}, \url{preprints/A-cubic-lower-bound-for-border-determinantal-complexity-of-the-permanent-September-24-2026}, commit \texttt{adc7f1241b42}.
\bibitem{Guo} J.~Guo, \emph{An $n^2\log\log n$ lower bound for permanent
circuits with valid division}, arXiv:2609.29568, v2, October 8, 2026.
\bibitem{Xie} Y.~Xie, \emph{Smooth slices and superquadratic border
determinantal complexity of the permanent}, Zenodo, record 23108131,
October 2, 2026.
\end{thebibliography}

\medskip\noindent\emph{Developed with AI assistance (Anthropic Claude, Opus 5.5)
under [AUTHOR NAME]'s direction; independent adversarial reviews were run by
separate AI sessions and are not human peer review.}
\end{document}
